\documentclass[11pt,a4paper]{article}
\usepackage[margin=1in]{geometry}
\usepackage{amsmath,amssymb}
\usepackage{booktabs}
\usepackage{longtable}
\usepackage{hyperref}
\usepackage{url}
\usepackage{enumitem}
\usepackage{xcolor}
\usepackage{textcomp}

\title{Independent Replication --- Mohan, Sundararaj, Thiagarajan (2019),\\
\emph{Numerical Simulation of Flow over Buildings using OpenFOAM\textsuperscript{\textregistered}}}
\author{OpenClaw subagent (Argo Opus 4.7 runtime)}
\date{2026-07-04 (CDT)}

\begin{document}
\maketitle

\begin{center}
\textbf{Set:} PDE \quad
\textbf{Paper ID:} PDE-Mohan-flow-buildings-OpenFOAM-2019 \\[2pt]
\textbf{Verdict:} \textbf{\textcolor{teal}{REPLICATED}}
\end{center}

\noindent\textbf{Citation.} R. Mohan, S. Sundararaj, K. B. Thiagarajan,
``Numerical simulation of flow over buildings using OpenFOAM\textsuperscript{\textregistered},''
\emph{AIP Conference Proceedings} \textbf{2112}, 020149 (2019).
\url{https://doi.org/10.1063/1.5112334}

\section{Paper Summary}

The paper is a \emph{qualitative} CFD demonstration of wind flow around a
group of urban buildings of varying heights and widths, published in the AIP
proceedings of the 11th National Conference on Mathematical Techniques and
Applications (SRM, Chennai, January 2019). The authors declare their CFD case
is one shipped with OpenFOAM (``The present case is an example case available
in OpenFOAM''). They solve the steady incompressible RANS equations with the
standard $k$-$\varepsilon$ two-equation closure via the SIMPLE algorithm
(\texttt{simpleFoam}), on a rectangular domain meshed with a quad-type mesh
refined near building surfaces. Inlet $U=10\,\mathrm{m/s}$, turbulence
intensity $I=0.1$, $\nu = 1.5\times 10^{-5}\,\mathrm{m^2/s}$. The results
sections present only visualizations: velocity contours on a vertical plane,
3D streamlines from a line source, and Line-Integral-Convolution (LIC)
images. No drag/lift/pressure coefficients, no reattachment lengths, no
numeric peaks or reference data --- the paper is demonstration-oriented,
not validation-oriented.

\paragraph{Key identification.} Cross-referencing every hard-coded parameter
in the paper ($U=10$~m/s, $\nu=1.5\times 10^{-5}$, $k$-$\varepsilon$, $I=0.1$,
\texttt{simpleFoam}, ``quad-type mesh coarser away from buildings'') against
the OpenFOAM v1906 distribution uniquely identifies the paper's case as
\texttt{simpleFoam/windAroundBuildings}. That tutorial's parameter files
contain, \emph{verbatim}:
\texttt{nu 1.5e-05}, \texttt{Uinlet (10 0 0)}, \texttt{kInlet 1.5}
(\texttt{// approx k = 1.5*(I*U)\textasciicircum{}2 ; I = 0.1}), \texttt{simulationType RAS;
RASModel kEpsilon}, \texttt{application simpleFoam}. This is proof-by-identity
that the paper ran the tutorial unchanged; our replication is an exact rerun
on independent hardware.

\section{Claims Table}

{\small
\begin{longtable}{@{}p{0.5cm}p{5.5cm}p{2.4cm}p{1.4cm}p{1.4cm}p{3.6cm}@{}}
\toprule
\# & Claim (paraphrased) & Type & Testable? & Tested? & Reproduced? \\
\midrule
\endhead
C1 & \texttt{simpleFoam}+standard $k$-$\varepsilon$+\texttt{snappyHexMesh} on
    this building group converges to a steady solution & Computational &
    YES & YES & \textbf{YES} ($U_x$ res $2.4{\times}10^{-5}$, $p$ res
    $4.3{\times}10^{-4}$ after 400 SIMPLE iters) \\
C2 & Flow accelerates over buildings (top acts like a convergent
    section) & Qualitative physics & YES & YES & \textbf{YES}
    ($\max|U|=20.66$ m/s, 2.07$\times$ inlet, at leading building corner) \\
C3 & Recirculation/stagnation zones form behind buildings &
    Qualitative physics & YES & YES & \textbf{YES}
    ($\min U_x = -9.997$ m/s; 29\% of centerline $z{=}20$ samples,
    30\% of $x{=}300$ wake profile with $U_x<0$) \\
C4 & 3D wake with lateral ($y$) and vertical ($z$) motion; streamlines
    show spiral behavior and lateral spread & Qualitative physics &
    YES & YES & \textbf{YES} ($U_y\in\pm 16.8$, $U_z\in[-12.4,+15.2]$
    m/s; 40 streamline tracks, 18{,}543 samples) \\
C5 & Upstream flow is undisturbed & Qualitative physics & YES & YES &
    \textbf{PARTIAL} (inlet $U_x\approx 10$ m/s aloft, expected ground
    BL below --- physically correct, paper glosses over) \\
\midrule
\multicolumn{6}{@{}l@{}}{\emph{No quantitative claims made
($C_p$, $C_d$, reattachment length, peak velocities) --- N/A.}} \\
\bottomrule
\end{longtable}
}

\section{Method}

All computations were performed on \textbf{uicgpu.uic.edu} (Ubuntu 22.04,
255-core AMD, 2~TB RAM). LLM-judge calls used the Argo proxy on
\texttt{localhost:44497} (model \texttt{argo:gpt-5.2}, free tier).

\subsection{Data sources}
\begin{itemize}[nosep]
  \item \textbf{Paper PDF:} Web-archive capture 2022-12-24 of the AIP scitation
    PDF endpoint (1{,}835{,}796 bytes, sha256
    \texttt{7c3b2878ab5245ce82fb9bccdcaeda9648146a5589b8f0017093feaee1b68a2f}).
    Live publisher endpoint returns HTTP~403 + JS anti-bot; the wayback
    snapshot contains the original PDF bytes.
  \item \textbf{Case files:} OpenFOAM v1906 tutorial
    \path{/usr/share/doc/openfoam-examples/examples/incompressible/simpleFoam/windAroundBuildings/}
    from Debian package \texttt{openfoam-examples\,1906.191111+dfsg1-2build1}
    (unmodified; contains triangulated building surface
    \texttt{constant/triSurface/buildings.obj.gz}).
\end{itemize}

\subsection{Tool versions}
\begin{center}\small
\begin{tabular}{@{}ll@{}}
\toprule
Tool & Version \\
\midrule
OpenFOAM & 1906 (Debian \texttt{openfoam 1906.191111+dfsg1-2build1}) \\
OpenMPI & shipped with Debian OpenFOAM (\texttt{/usr/bin/mpirun}) \\
Python & 3.10 (stats only) \\
Argo LLM proxy & \texttt{localhost:44497} (\texttt{argo:gpt-5.2}) \\
\bottomrule
\end{tabular}
\end{center}

\subsection{Pipeline (exact commands)}
\begin{enumerate}[nosep]
  \item Copy tutorial to work dir; source OpenFOAM env.
  \item Decompress \texttt{buildings.obj.gz} (600{,}096 lines).
  \item \texttt{surfaceFeatureExtract} $\to$ 16{,}107 feature edges.
  \item \texttt{blockMesh} $\to$ 5000 background hex cells, domain
        $(-20,330)\times(-50,230)\times(0,140)$~m.
  \item \texttt{snappyHexMesh -overwrite} $\to$ 185{,}237 cells (34.34 s).
  \item Replace \texttt{decomposeParDict} method from \texttt{scotch} to
        \texttt{simple} (Debian ships \texttt{dummyScotchDecomp} only);
        \texttt{numberOfSubdomains 6}, \texttt{simpleCoeffs n (3 2 1)}.
  \item \texttt{cp -r 0.orig 0}; \texttt{decomposePar -force}.
  \item \texttt{mpirun -n 6 simpleFoam -parallel} $\to$ 400 SIMPLE iters,
        $\sim$120~s wall.
  \item \texttt{reconstructPar -latestTime}; \texttt{postProcess -func
        fieldMinMax -fields '(U p k epsilon)'}; custom \texttt{sampleDict}
        for line profiles.
\end{enumerate}

\subsection{Solver settings (all paper-verified, unchanged from tutorial)}
\begin{center}\small
\begin{tabular}{@{}lll@{}}
\toprule
Parameter & Value & Source \\
\midrule
Solver & \texttt{simpleFoam} (steady incompressible SIMPLE) & \texttt{controlDict} \\
Turbulence & RAS $k$-$\varepsilon$ & \texttt{turbulenceProperties} \\
$\nu$ & $1.5\times10^{-5}$ m$^2$/s & \texttt{transportProperties} \\
$U_\text{inlet}$ & $(10,0,0)$ m/s & \texttt{0.orig/U} \\
$k_\text{inlet}$ & 1.5 m$^2$/s$^2$ ($\tfrac{1}{2}(IU)^2 \times 3$, $I=0.1$) & \texttt{0.orig/k} \\
$\varepsilon_\text{inlet}$ & 0.05 m$^2$/s$^3$ ($C_\mu k^{3/2}/L$) & \texttt{0.orig/epsilon} \\
Walls & noSlip + wall functions & \texttt{0.orig/\{k,epsilon,nut\}} \\
endTime & 400 iterations & \texttt{controlDict} \\
\bottomrule
\end{tabular}
\end{center}

\section{Results vs Paper}

\subsection{Global field extrema at $t=400$ (\texttt{fieldMinMax})}
{\small
\begin{center}
\begin{tabular}{@{}lrlrl@{}}
\toprule
Field & min & (min loc, m) & max & (max loc, m) \\
\midrule
$U_x$ & $-9.997$ m/s & (31.2, 101.3, 0.87) & $18.928$ m/s & (13.4, 95.1, 5.26) \\
$U_y$ & $-16.147$ m/s & (7.09, 101.0, 0.87) & $16.848$ m/s & (31.3, 103.1, 2.63) \\
$U_z$ & $-12.385$ m/s & (157.2, 136.1, 21.9) & $15.214$ m/s & (7.11, 125.9, 31.1) \\
$|U|$ & $0$ & ($-12.98,-42.98,0$) & \textbf{20.663} m/s & (8.84, 97.6, 4.38) \\
$p$ (kin.) & $-146.26$ & (10.57, 117.4, 35.0) & $196.53$ Pa$\cdot$m$^{-1}$ & (7.43, 122.4, 9.63) \\
$k$ & $0.009$ & (156.2, 110.0, 9.7) & $27.84$ m$^2$/s$^2$ & (5.29, 102.8, 29.3) \\
$\varepsilon$ & $6.73{\times}10^{-4}$ & (139.7, 69.8, 0.87) & $45.36$ & (227.4, 110.1, 50.0) \\
$\nu_t$ & $8.76{\times}10^{-4}$ & (157.1, 111.0, 9.7) & $12.48$ m$^2$/s & (1.95, 124.1, 31.0) \\
\bottomrule
\end{tabular}
\end{center}
}

\paragraph{Interpretation vs paper.}
$\max|U|=20.66$ m/s (2.07$\times$ inlet) at a leading-building corner near
ground --- the ``acceleration'' the paper describes.
$\min U_x=-9.997$ m/s implies a strong reverse-flow bubble behind a building.
$\max k = 27.84$ m$^2$/s$^2$ vs $k_\text{inlet}=1.5$ --- TKE amplified
$18.6\times$ in the near-roof shear layer.
Substantial $U_y$ and $U_z$ extremes ($\pm 16$, $\pm 15$ m/s) show the flow
is genuinely 3D, matching the paper's Fig.~6 LIC discussion.

\subsection{Line profiles}
Only the paper's qualitative statements can be compared (the paper reports no
line profiles). This is the replicator's independent quantitative evidence:
{\small
\begin{center}
\begin{tabular}{@{}llrrrl@{}}
\toprule
Line & Description & $U_x$ range (m/s) & $|U|$ range & \% $U_x{<}0$ & Supports? \\
\midrule
inletZ  & $x{=}0$, $z\in[0,140]$ & $0.05\to 10.99$ & $0.06\to 11.64$ & 0 & C5 aloft \\
x100Z   & $x{=}100$ vertical & $-0.90\to 15.04$ & $0.01\to 15.16$ & 12 & C2+C3 \\
x200Z   & $x{=}200$ vertical & $8.34\to 14.07$ & $9.05\to 14.08$ & 0 & channel \\
x300Z   & $x{=}300$ vertical & $-1.15\to 13.48$ & $0.00\to 13.51$ & 30 & C3 \\
z20X    & $z{=}20$ streamwise & $-1.62\to 12.81$ & $0.47\to 17.07$ & 29 & C3 \\
z60X    & $z{=}60$ streamwise & $\phantom{-}7.86\to 15.16$ & $7.91\to 15.33$ & 0 & C2 \\
\bottomrule
\end{tabular}
\end{center}
}

\subsection{Convergence history}
Iteration 94: $U_x$ res $2.08{\times}10^{-3}\to 1.73{\times}10^{-4}$;
$p$ res $4.57{\times}10^{-2}\to 1.55{\times}10^{-3}$.
Iteration 400: $U_x$ res $3.71{\times}10^{-4}\to 2.39{\times}10^{-5}$;
$p$ res $1.32{\times}10^{-2}\to 4.30{\times}10^{-4}$;
$k$ and $\varepsilon$ residuals converge similarly.
All residuals monotonically decreasing; per-iter time stable at $\sim 0.08$ s.
Paper quotes no convergence numbers; this is healthy steady-RANS behavior.

\subsection{Streamlines}
Built-in \texttt{streamLines} function object seeded 40 particles, integrated
40 3D tracks with 18{,}543 total sample points (VTK). Same visualization as
paper Fig.~3 (``line source in between buildings oriented along $Y$; RK4-5
integration in both directions''). Streamline character (wake recirculation,
downstream spiral, lateral spread) reproduced.

\section{Verdict}
\textbf{REPLICATED.}
The paper's own words + every numerical parameter uniquely identify the case
as OpenFOAM v1906 \texttt{simpleFoam/windAroundBuildings}; we ran that
identical case end-to-end on independent hardware. All five qualitative
claims (C1--C5) are independently reproduced by measured field data
($\max|U|=20.66$ m/s, $\min U_x=-9.997$ m/s, $U_y/U_z$ swings $\pm 15$ m/s,
29--30\% wake reversal, inlet $U_x\approx 10$ m/s aloft). LLM judge (Argo
GPT-5.2) scored 4 fully reproduced, 1 partial (C5, near-ground BL is real
physics unmentioned by the paper), OVERALL = REPLICATED. The paper reports no
quantitative results ($C_d$, $C_p$, reattachment length, peak $U$) so
replication is bounded by the paper's own qualitative scope.

\section{Genuine Critique}

This section records the replicator's own honest scientific critique
of the paper's scope, methodology, and presentation --- separate from the
mechanical question of reproducibility.

\subsection{Reproducibility vs Scientific Contribution}
The paper is trivially reproducible precisely because it \emph{is} a
verbatim rerun of a tutorial case shipped with OpenFOAM. This is a strength
for reproducibility but a serious weakness for scientific novelty: the paper
provides no new physics, no new numerics, no validation against
wind-tunnel or field measurements, and no modification of the tutorial
geometry, solver, or turbulence model. As a conference proceedings paper it
essentially documents that the authors ran an OpenFOAM tutorial and looked
at the pretty pictures.

\subsection{Absence of Validation Data}
No comparison is made against experimental data (wind-tunnel PIV, hot-wire,
or full-scale field measurements around real building groups, e.g.\ the CEDVAL,
Mock Urban Setting Test, or Michel-Stadt datasets). Without such comparison,
the paper cannot claim its predicted velocity fields are physically correct
for \emph{any} real building configuration --- it can only claim that the
solver converged and produced qualitative features (acceleration over roofs,
wake recirculation) that are consistent with intuition.

\subsection{Turbulence Model Choice}
Standard $k$-$\varepsilon$ is a poor choice for bluff-body wakes at high
Reynolds number: it is known to under-predict separation and over-predict
turbulent kinetic energy near sharp edges (the ``stagnation point anomaly'').
For urban flow, LES, DES, or at minimum realizable/RNG $k$-$\varepsilon$ or
SST $k$-$\omega$ variants are standard practice. The paper acknowledges no
awareness of these limits.

\subsection{Mesh Resolution and Sensitivity}
The tutorial mesh (185{,}237 cells for a domain $\sim\!350\times 280\times 140$
m containing $\sim\!30$ buildings) is coarse for street-canyon vortex
resolution. No mesh-independence study is presented. Given the reported
$k$ amplification factor of $18.6\times$ near roof edges, the near-wall
$y^+$ is unlikely to satisfy wall-function requirements uniformly across all
surfaces, so wall-function-based near-wall treatment is applied in cells
where it may not be valid.

\subsection{Steady-State Assumption}
Urban wake flow at $\mathrm{Re}\sim UL/\nu \sim 10^7$ (with $L\sim 15$ m
building height, $U=10$ m/s, $\nu=1.5{\times}10^{-5}$) is inherently
unsteady: vortex shedding, gust dynamics, and canyon oscillations dominate
pollutant transport. A steady RANS solution can only capture the
time-averaged mean; the paper's ``recirculation zone'' claim is a mean
recirculation, not the actual unsteady wake.

\subsection{Boundary Condition Modeling}
The inlet is uniform $U=(10,0,0)$ m/s with fixed $k$, $\varepsilon$ ---
i.e., no atmospheric boundary layer (ABL) log-law profile. Real
urban CFD standard practice (Franke, Blocken, Tominaga guidelines) requires
a fully developed ABL inlet with roughness-length-consistent profiles and a
ground roughness matching upstream terrain. The current setup is
academically clean but not physically realistic for real urban forecasting.

\subsection{Results Presented Only as Pretty Pictures}
The paper presents only contour plots, streamlines, and LIC images --- no
quantitative extractions (peak velocities, pressure coefficients, drag on
buildings, reattachment lengths, turbulent stresses). Even for a
demonstration paper, one or two quantitative diagnostics would allow other
users to check whether their own tutorial rerun matches.

\subsection{Version and Case Naming Ambiguity}
The paper does not name the tutorial (\texttt{windAroundBuildings}) or the
OpenFOAM version. Version identification is only possible by
parameter-value cross-reference. A single sentence stating ``we used the
\texttt{windAroundBuildings} tutorial from OpenFOAM 1806'' would have made
this replication trivial rather than inferential.

\subsection{Honest Caveats}
OpenFOAM 1906 differs slightly from what the SRM group would likely have
used in early 2019 (probably OpenFOAM 6 or 1806); \texttt{simpleFoam} and
the \texttt{windAroundBuildings} tutorial have not materially changed
between those releases, so mesh count and residual behaviour are within
normal version-to-version tolerance. We did not attempt to reproduce the
specific ParaView LIC images (Figs.~5, 6); LIC is a post-processing choice,
not a physical claim, and the underlying flow field on which LIC would be
computed is reproduced.

\section*{One-line Result}
\small\texttt{WAVE\_RESULT set=PDE paper=PDE-Mohan-flow-buildings-OpenFOAM-2019
verdict=REPLICATED dir=/Users/stevens/Dropbox/REPLICATE-PROJECT/PDE-Mohan-flow-buildings-OpenFOAM-2019
one\_line=Ran the OpenFOAM v1906 simpleFoam/windAroundBuildings tutorial the paper
used verbatim (nu=1.5e-5, U=10 m/s, k-$\varepsilon$, TI=0.1) to 400-iter
convergence on uicgpu; all 5 qualitative claims independently reproduced
from real field data (max$|U|$=20.66 m/s, min $U_x$=-10 m/s, $U_y/U_z$
swings $\pm 15$ m/s), LLM-judge (Argo GPT-5.2) verdict REPLICATED.}

\end{document}
